\begin{definition}[Native bond supporting a charge]
    \label{def:peps_torus_edge_physical_charge_measurement}
    \lean{TNLean.PEPS.torusChargeMeasurementEdge}
    \leanok
    \uses{def:peps_regular_edge_charge_contraction}
    On a finite simple torus, choose either the right edge or the up edge of a
    native vertex $v$. Its ordered endpoints form the measured two-site region.
    The diagonal insertion is $\chi(pk)$ on this actual shared bond, including
    bonds that cross either coordinate seam.
\end{definition}

\begin{theorem}[Nonzero native two-spin charged columns]
    \label{thm:peps_torus_edge_charge_nonzero}
    \lean{TNLean.PEPS.IsGIsometric.openRegionWeight_torusEdgeCharacterSite_ne_zero}
    \leanok
    \uses{def:peps_torus_edge_physical_charge_measurement,
        thm:peps_regular_two_site_charge_nonzero,lem:peps_regular_torus_incident_isometry}
    Suppose both torus periods are at least three and the original four-leg tensor
    is regular $G$-isometric. For every actual irreducible character $\chi$, every
    $p\in G$, and every native boundary configuration $\theta$, the charged
    open column on the two endpoints of a chosen right or up edge is nonzero.
    This asserts no nonvanishing under an arbitrary global exterior contraction.
\end{theorem}
\begin{proof}\leanok
    At each endpoint a perpendicular incident edge differs from the shared edge.
    Thus both sites have remaining boundary legs. Native incident-edge coordinates
    preserve regular $G$-isometry, and hence $G$-injectivity. The nonzero charged
    column theorem applies; the actual contraction identity and physical endpoint
    bijection preserve nonvanishing.
\end{proof}

\begin{theorem}[Native complete two-spin charge detection]
    \label{thm:peps_torus_edge_physical_charge_measurement}
    \lean{TNLean.PEPS.IsGIsometric.exists_torusEdgePhysicalAllChargeMeasurement}
    \leanok
    \uses{thm:peps_regular_edge_physical_charge_measurement,
        lem:peps_regular_torus_incident_isometry,thm:peps_physical_cut_column_action}
    Suppose both torus periods are at least three and the original four-leg tensor
    is regular $G$-isometric. At every native right or up edge there is one complete
    projective measurement $Q$ on its two original physical spins such that
    \begin{align}
        Q_r\Psi^{\chi,p}_{R_e,\theta}
         & =\mathbf1_{r=\chi}\Psi^{\chi,p}_{R_e,\theta}, \\
        Q_r C^{\chi,p}_{R_e}
         & =\mathbf1_{r=\chi}C^{\chi,p}_{R_e}.
    \end{align}
    Here $C^{\chi,p}_{R_e}$ is the actual coefficient matrix of the global cut.
    The same measurement precedes every irreducible charge, internal parameter,
    boundary label, and arbitrary exterior site family agreeing with the original
    tensors at the measured pair. Both seam orientations are included by using
    the actual ordered edge. This is the finite-torus specialization of
    \cite[Detection of chargeons]{Schuch2010PEPS}, lines 2464--2486; it includes
    no parent-Hamiltonian assertion.
\end{theorem}
\begin{proof}\leanok
    Use the complete physical pair measurement in native incident-edge
    coordinates. Its action holds on every original open-region column. The
    actual region-complement contraction expresses the cut matrix as the open
    matrix times the transpose of the complementary matrix, so the same
    indicator action holds on the cut for arbitrary exterior tensors.
\end{proof}
